\begin{abstract}
In most cases, when both of them are loaded into the accelerator, the problems are likely to emerge sooner than anticipated. For instance, memory on the power infrastructure proved to be a bigger bottleneck than either latency or power delivery.
The volatility consistently observed in benchmarks had little to do with computing performance. As the workset moved along the hierarchy, interactions between GPU video RAM (VRAM), host dynamic RAM (DRAM), and NVMe (Non-Volatile Memory Express) became problematic. After this balancing was disrupted, the system started to behave in an unpredictable way. The latency became really high, but the profiler could not give us an adequate answer about what caused this instability.
This paper's contribution lies in formally analyzing these interactions. End-to-end latency consists of four elements: computing, localization, inter-layer transfer, and synchronization. Decomposition yields two ratios: the capacity ratio $\alpha$ and the input/output (I/O) pressure ratio $\beta$. This ratio represents the point at which memory throttling is actually effective, i.e., the point at which the system has reached its bandwidth limit.
Then, we use these ratios to construct a RAMSE. The RAMSE is small in size (runtime control plane + 3 modules) and controls VRAM, DRAM, and NVMe according to a system-aware policy. The results are as follows: Model load time decreased by up to 60\%, inference latency by 35\%, maximum VRAM usage decreased by 15\%, and throughput per watt increased by 18.5\%. Stable performance was maintained even after 72 hours of stress testing, and the yield outside the SLA decreased by almost tenfold.
Correct interpretation of the collected data leads to the conclusion that, in industrial inference, hierarchical memory is not a consistently difficult problem, but rather exhibits different characteristics depending on the system. The boundary values we derived provide a benchmark that system designers can specifically refer to when aiming for stable and energy-efficient inference within industrial information system stacks.
\end{abstract}

\begin{IEEEkeywords}
Industrial Informatics, Industrial AI, Edge Computing, Deterministic Latency, Hierarchical Memory Management, GPU Offloading
\end{IEEEkeywords}
